\documentclass[a4paper,11pt]{ltjsarticle}
\input{preamble}

\title{電場と電位:2つの見方をつなぐ(動作確認用の例)}
\author{}
\date{\today}

\begin{document}
\maketitle

\section{この単元で何が分かると嬉しいか}

電荷のまわりの空間を「向きと大きさを持つ場」として見るのが\textbf{電場}、
「高さのような1つの数」として見るのが\textbf{電位}である。
同じ現象を2通りに書いているだけなので、片方から他方へ移れるようになると、
問題に合わせて楽な方で解けるようになる。

\section{電場の定義}

点電荷 $q_1$ が距離 $r$ 離れた電荷 $q_2$ に及ぼす力の大きさはクーロンの法則で与えられる。

\begin{fact}{クーロンの法則}
  \[
    F = \frac{1}{4\pi\eps_0}\,\frac{q_1 q_2}{r^2}
  \]
  同符号なら反発、異符号なら引き合う。$\eps_0$ は真空の誘電率。
\end{fact}

試しに置いた電荷 $q$ が受ける力を $\vect{F}$ とするとき、電荷の大きさで割った量を電場とする。
\begin{equation}
  \vect{E} = \frac{\vect{F}}{q}
  \label{eq:E-def}
\end{equation}
式~\eqref{eq:E-def} の右辺は $q$ に依らないので、電場は「その場所が持つ性質」になる。
原点にある点電荷 $Q$ がつくる電場の大きさは $E = Q/(4\pi\eps_0 r^2)$ である。

\begin{point}
  電場は「置いた電荷1\,\unit{C}あたりの力」。単位は \unit{N/C}(= \unit{V/m})。
\end{point}

\section{電位との関係}

点電荷のまわりの電位は、無限遠を基準として
\begin{equation}
  \phi(r) = \frac{Q}{4\pi\eps_0 r}
\end{equation}
となる。電場は電位の傾き(の符号を変えたもの)である。
\begin{equation}
  \vect{E} = -\grad\phi, \qquad
  E_r = -\dv{\phi}{r} = \frac{Q}{4\pi\eps_0 r^2}
\end{equation}
右端が先の $E$ と一致することを確かめよ、というのが次の確認問題である。

\begin{caution}
  電位は向きを持たない数(スカラー)。電場の向きは「電位が下がる向き」なので、
  符号のマイナスを落とすと向きが逆になる。
\end{caution}

\begin{question}
  $\phi(r) = Q/(4\pi\eps_0 r)$ を $r$ で微分して、$E_r = Q/(4\pi\eps_0 r^2)$ を導け。
\end{question}

\end{document}
